\chapter{Graf jako przestrzeń topologiczna}
\label{ch:tg-topology}

Dotychczas traktowaliśmy graf jako zbiór wierzchołków i krawędzi.
Od teraz każdą krawędź zastępujemy rzeczywistym odcinkiem.
Tak powstała przestrzeń jest kompleksem CW o wymiarze
co najwyżej jeden.
Dla topologii ważne jest, że nawet jeśli narysujemy graf
na płaszczyźnie z przecięciami, nie wprowadzamy tam
nowych wierzchołków bez wyraźnej identyfikacji.

\section{Kompleks CW, charakterystyka Eulera i zwijanie drzewa}
\label{sec:tg-collapse-tree}

Przypomnijmy konstrukcję realizacji geometrycznej $|G|$
z rozdziału \ref{ch:tg-podstawy}: dla każdego $v\in V$ bierzemy
$0$-komórkę, a dla każdej krawędzi $e$ odcinek z końcami
doklejonymi do odpowiednich wierzchołków. Krawędzie
równoległe i pętle również dają poprawne CW-kompleksy,
gdy zachowamy osobne odcinki i odrębnie określimy odwzorowania
przyklejania ich końców.

\begin{definicja}[Charakterystyka Eulera]
Dla skończonego kompleksu CW $X$ oznaczamy przez $c_k(X)$
liczbę jego otwartych $k$-komórek i definiujemy
$\chi(X)=\sum_{k\ge0}(-1)^k c_k(X)$.
Dla realizacji skończonego grafu $G$ otrzymujemy
$\chi(|G|)=|V(G)|-|E(G)|$.
\end{definicja}

\begin{lemat}[Zwinięcie kontraktowalnego podkompleksu]
\label{lem:tg-collapse-CW}
Niech $X$ będzie kompleksem CW, a $A\subseteq X$
niepustym, kontraktowalnym podkompleksem CW.
Wtedy naturalne odwzorowanie ilorazowe $q:X\to X/A$
jest równoważnością homotopijną.
\end{lemat}

\begin{proof}
Ponieważ $(X,A)$ jest parą CW, ma własność przedłużania
homotopii: każdą homotopię $A\times[0,1]\to X$,
rozpoczynającą się od włożenia $A\hookrightarrow X$,
można przedłużyć do homotopii $H:X\times[0,1]\to X$
z $H_0=\operatorname{id}_X$. Wybieramy kontrakcję
$A$ do punktu $a_0\in A$, tak aby $H_t(A)\subseteq A$
i $H_1(A)=\{a_0\}$. Wtedy $H_1$ jest stałe na $A$,
więc faktoryzuje się przez iloraz jako $H_1=g\circ q$
dla pewnego ciągłego $g:X/A\to X$.

Z jednej strony $gq=H_1\simeq\operatorname{id}_X$.
Z drugiej strony $qH_t$ jest stałe na $A$ dla każdego $t$,
więc daje homotopię na ilorazie od
$\operatorname{id}_{X/A}$ do $qg$. Ciągłość tej homotopii
wynika z tego, że $q\times\operatorname{id}_{[0,1]}$
jest odwzorowaniem ilorazowym (iloczyn odwzorowania
ilorazowego z identycznością na lokalnie zwartej
przestrzeni Hausdorffa zachowuje tę własność).
Stąd również $qg\simeq\operatorname{id}_{X/A}$.
Obie równości homotopijne dowodzą tezy.
\end{proof}

\begin{uwaga}
W dowodzie użyto standardowej własności par CW.
Zwijanie dowolnego kontraktowalnego podzbioru przestrzeni
bez założenia o podkompleksie lub odpowiednim warunku
kofibracji nie musi dawać równoważności homotopijnej. Warunek ten jest istotny.
\end{uwaga}

\newpage
\begin{twierdzenie}[Graf jest homotopijnie równoważny bukietowi okręgów]
\label{thm:tg-wedge}
Niech $G$ będzie skończonym spójnym grafem nieskierowanym
z $n\ge1$ wierzchołkami i $m$ krawędziami
(dopuszczamy pętle i krawędzie równoległe).
Dla dowolnego drzewa rozpinającego $T\subseteq G$
przestrzeń $|G|$ jest homotopijnie równoważna bukietowi
$\bigvee_{j=1}^{m-n+1}S^1$.
Przy zerowej liczbie okręgów bukiet oznacza pojedynczy punkt.
\end{twierdzenie}

\begin{proof}
Drzewo $T$ jest kontraktowalne. Jeżeli ma co najmniej
dwa wierzchołki, wybieramy liść $v$ i jego jedyną
krawędź, parametryzowaną przez $[0,1]$ od sąsiada do $v$.
Na tej krawędzi homotopia ma postać $s\mapsto(1-t)s$
dla $t\in[0,1]$, a na reszcie drzewa jest identycznością.
Obie definicje zgadzają się na wspólnym końcu $s=0$,
więc dają ciągłą retrakcję deformacyjną na drzewo
bez liścia i jego krawędzi. Ponieważ drzewo
jest skończone, po skończonej liczbie deformacji
pozostaje pojedynczy wierzchołek.

Z lematu \ref{lem:tg-collapse-CW} iloraz
$q:|G|\to|G|/|T|$ jest równoważnością homotopijną.
W ilorazie wszystkie wierzchołki drzewa, czyli wszystkie
wierzchołki $G$, stają się jednym punktem.
Każda krawędź spoza $T$ pozostaje odcinkiem, którego
oba końce zostały ze sobą utożsamione, a więc tworzy
okrąg. Różne okręgi mają wspólny tylko ten punkt.
Ich liczba wynosi $m-(n-1)=m-n+1$.
\end{proof}

\begin{center}
\begin{tikzpicture}[scale=.96,line cap=round,>=Stealth]
 \node[tgBlueVertex] (a) at (-3.9,.8) {$a$};
 \node[tgBlueVertex] (b) at (-2.1,1.25) {$b$};
 \node[tgBlueVertex] (c) at (-1.8,-.65) {$c$};
 \node[tgBlueVertex] (d) at (-3.7,-.8) {$d$};
 \draw[tgTree] (a)--(b)--(c)--(d);
 \draw[tgAccent] (d)--(a) (a)--(c);
 \draw[tgTeal,-{Stealth[length=2.6mm]},line width=1.6pt] (-.85,.15)--(.25,.15);
 \node[font=\footnotesize,text=tgNavy] at (-.25,.51) {zwiń $T$};
 \node[tgBlueVertex] (o) at (2.0,.15) {$*$};
 \draw[tgAccent,line width=2pt] (o) .. controls (1.2,1.9) and (3.6,1.9) .. (o);
 \draw[tgAccent,line width=2pt] (o) .. controls (1.2,-1.4) and (3.6,-1.4) .. (o);
 \node[font=\footnotesize,text=tgNavy] at (1.95,-1.78) {bukiet dwóch okręgów};
\end{tikzpicture}

{\footnotesize Zwijamy zielone drzewo. Dwie czerwone krawędzie
pozostają generatorami dwóch niezależnych okręgów.}
\end{center}

\section{Grupa fundamentalna grafu}
\label{sec:tg-pi1}

\begin{definicja}[Grupa wolna]
Dla zbioru $S$ przez $F(S)$ oznaczamy grupę, której elementami
są zredukowane słowa z liter $s$ i $s^{-1}$, $s\in S$,
z mnożeniem przez konkatenację i redukcję sąsiadujących par
$s s^{-1}$ lub $s^{-1}s$. Dla $|S|=r<\infty$ piszemy $F_r$.
Grupą $F_0$ jest grupa trywialna.
\end{definicja}

\begin{twierdzenie}[Grupa fundamentalna skończonego grafu]
\label{thm:tg-free-pi1}
Niech $G$ będzie skończonym spójnym grafem nieskierowanym
o $n\geq1$ wierzchołkach i $m$ krawędziach
(dopuszczamy pętle i krawędzie równoległe),
a $x_0\in|G|$ dowolnym punktem bazowym. Jeżeli
$\beta=m-n+1$, to $\pi_1(|G|,x_0)\cong F_\beta$.
W szczególności $|G|$ jest po prostu spójna wtedy i tylko
wtedy, gdy $G$ jest drzewem.
\end{twierdzenie}

\begin{proof}
Wybierzmy drzewo rozpinające $T$ i wierzchołek bazowy
$r\in T$. Poprzednie twierdzenie daje równoważność
homotopijną $|G|\simeq\bigvee_\beta S^1$.
Dla $\beta=0$ bukiet jest punktem i jego grupa
fundamentalna jest trywialna. Dla $\beta\geq1$
grupa fundamentalna bukietu $\beta$ okręgów jest wolna
na $\beta$ generatorach. Przypomnijmy uzasadnienie
przez twierdzenie van Kampena: dla każdego okręgu
wybieramy zbiór otwarty złożony z całego tego okręgu
i małych łuków wokół punktu sklejenia na pozostałych
okręgach. Wybieramy te same małe otwarte łuki zawierające
punkt sklejenia na każdym okręgu. Każde otoczenie
deformacyjnie ściąga się na swój pełny okrąg
przez kontrakcję pozostałych łuków. Przecięcie
kolejnego otoczenia z sumą poprzednich jest sumą
wspomnianych małych łuków, kontraktowalną do punktu
sklejenia. Wszystkie zbiory i te przecięcia są
spójne drogowo. Iteracja twierdzenia van Kampena
daje wolny iloczyn $\mathbb Z*\cdots*\mathbb Z=F_\beta$.

Generatorowi odpowiada pętla: od $r$ idziemy po jedynej
ścieżce drzewowej do początku wybranej krawędzi
$e\notin T$, przechodzimy $e$ i wracamy po drzewie do $r$.
Zmiana punktu bazowego wewnątrz spójnego $|G|$ nie zmienia
typu izomorfizmu grupy fundamentalnej.

Jeśli $\beta=0$, graf nie ma cykli i jest drzewem.
Gdy $\beta\ge1$, grupa $F_\beta$ zawiera nietrywialną
pętlę, a więc $|G|$ nie jest po prostu spójna.
\end{proof}

\begin{przyklad}[Trójkąt i graf theta]
Dla trójkąta $K_3$ mamy $n=m=3$ i $\beta=1$.
Zatem $|K_3|\simeq S^1$ oraz
$\pi_1(|K_3|)\cong\mathbb Z$.

Graf theta składa się z dwóch wierzchołków połączonych
trzema różnymi krawędziami (jest to multigraf bez pętli).
Ma $n=2$, $m=3$, więc $\beta=2$.
Jego realizacja jest homotopijnie równoważna bukietowi
dwóch okręgów, a grupa fundamentalna jest izomorficzna
z wolną grupą $F_2$, która jest nieprzemienna.
To pokazuje, że przestrzenie jednowymiarowe mogą mieć
nietrywialne, nawet nieprzemienne grupy fundamentalne.
\end{przyklad}

\begin{center}
\begin{tikzpicture}[scale=.98,line cap=round]
 \node[tgBlueVertex] (u) at (-1.75,0) {$u$};
 \node[tgBlueVertex] (v) at (1.75,0) {$v$};
 \draw[tgBlue,line width=2.1pt] (u) to[bend left=48] (v);
 \draw[tgTeal,line width=2.1pt] (u)--(v);
 \draw[tgCoral,line width=2.1pt] (u) to[bend right=48] (v);
 \node[text=tgNavy,font=\footnotesize] at (0,-1.22) {graf theta, $\pi_1\cong F_2$};
\end{tikzpicture}
\end{center}

\section{Nakrycie uniwersalne i drzewa}
\label{sec:tg-universal-cover}

Sytuacja staje się jeszcze ciekawsza, gdy zastąpimy skończony graf
jego nakryciem. Nakrycie jest lokalnie takie samo, ale może mieć
nieskończenie wiele wierzchołków i krawędzi.

\begin{definicja}[Nakrycie grafów]
Dla grafów $\widetilde G,G$ o realizacjach geometrycznych
odwzorowanie ciągłe i surjektywne
$p:|\widetilde G|\to|G|$ nazywamy
\emph{nakryciem grafowym}, jeżeli przenosi wierzchołki na
wierzchołki, każdą otwartą krawędź homeomorficznie na
otwartą krawędź, jest zgodne z odwzorowaniem końców
krawędzi oraz dla każdego wierzchołka
$\widetilde v$ indukuje bijekcję między zbiorami
\emph{incydencji} (końców krawędzi wychodzących z
$\widetilde v$) i $p(\widetilde v)$.
Dopuszczamy także grafy nieskończone, lokalnie skończone;
każda pętla wnosi do zbioru incydencji dwa końce.
\end{definicja}

Dla realizacji jednowymiarowych podany warunek jest dokładnie
lokalnym warunkiem nakryciowym: małe gwiazdowe otoczenie
każdego wierzchołka ma rozłączne kopie w przestrzeni nakrywającej.

\begin{twierdzenie}[Nakrycie uniwersalne realizacji grafu]
\label{thm:tg-universal}
Każdy spójny, skończony graf $G$ posiada spójne nakrycie
$p:|\widetilde G|\to|G|$, w którym $\widetilde G$ jest drzewem.
Nakrycie to jest uniwersalne.
\end{twierdzenie}

\begin{proof}[Konstrukcja ze zredukowanych ścieżek]
Wybierzmy wierzchołek $r\in G$ i orientację pomocniczą
krawędzi. Wierzchołkami $\widetilde G$ uczynimy wszystkie
\emph{zredukowane spacery krawędziowe} rozpoczynające się w $r$:
nie dopuszczamy dwóch kolejnych kroków po tej samej krawędzi
w przeciwnych kierunkach. Spacer pusty jest korzeniem $\widetilde r$.
Każdy taki spacer kończy się w pewnym wierzchołku $v\in G$.
Każdy niepusty zredukowany spacer $w$ określa dokładnie
jedną krawędź w $\widetilde G$, łączącą $w$ ze spacerem
powstałym przez skreślenie jego ostatniego kroku.
Tym samym nie liczymy tej samej krawędzi dwukrotnie.
Z punktu $w$ można wykonać każdy krok po incydentnej
krawędzi $e$ w kierunku od wierzchołka końcowego $w$:
jeśli jest to odwrotność ostatniego kroku,
przechodzimy do poprzednika, a w przeciwnym razie
do spaceru uzyskanego przez dopisanie tego kroku.
W przypadku pętli rozróżniamy oba kierunki jej przejścia.
Różne krawędzie równoległe w $G$ pozostają rozróżnione.

Rzutowanie $p$ przyporządkowuje spacerowi jego koniec w $G$,
a każdej nowej krawędzi odpowiadającą jej krawędź bazową.
Przy każdym wierzchołku $w$ i każdym możliwym kroku w $G$
istnieje dokładnie jedno odpowiednie przejście w $\widetilde G$.
Każdy wierzchołek grafu bazowego jest końcem pewnej
ścieżki z $r$, która jest zredukowanym spacerem.
Każdą krawędź można podnieść, wychodząc z dowolnego
przeciwobrazu jej początku, dzięki lokalnej bijekcji.
Zatem $p$ jest surjektywne także na wnętrzach krawędzi
i spełnia wszystkie warunki nakrycia.

Każdy niepusty zredukowany spacer ma jednego poprzednika:
otrzymujemy go przez usunięcie ostatniego kroku.
Każda inna krawędź wychodząca z tego wierzchołka wydłuża
spacer o jeden krok. Wobec tego każdy wierzchołek można
połączyć z $\widetilde r$ przez jednoznaczny ciąg poprzedników.
Cykl w $\widetilde G$ musiałby mieć wierzchołek największej
długości spaceru; miałby on na cyklu dwóch różnych sąsiadów
o mniejszej długości, mimo że istnieje tylko jeden poprzednik.
To niemożliwe. Zatem $\widetilde G$ jest spójnym drzewem.

W metryce długości (krawędzie mają długość jeden)
kontrakcję $H_u$ otrzymujemy, przesuwając $x$ ku korzeniowi
$\widetilde r$ do odległości $(1-u)d(\widetilde r,x)$. Drzewo jest lokalnie skończone, gdyż lokalna bijekcja
incydencji zachowuje skończone stopnie grafu bazowego.
Metryka wyznacza topologię realizacji: małe kule wokół
wierzchołka są gwiazdami odcinków, a skończona liczba
ramion pozwala wybrać wspólny dodatni promień w każdym
otoczeniu CW; we wnętrzu krawędzi mamy zwykłe przedziały.

Niech $a=d(\widetilde r,x)$, $b=d(\widetilde r,y)$,
a $c$ będzie długością wspólnego początkowego odcinka
jedynych dróg od korzenia do $x$ i $y$. Wówczas
$d(x,y)=a+b-2c$. Po przesunięciu punktów na odległości
$sa,sb$, gdzie $s=1-u$, wspólna część ich dróg
ma długość $\min\{c,sa,sb\}$. Jeżeli minimum wynosi $c$,
nowa odległość to $s(a+b)-2c\leq d(x,y)$.
W przeciwnym razie wynosi ona $s|a-b|\leq d(x,y)$,
gdyż $c\leq\min\{a,b\}$. Stąd
$d(H_u(x),H_u(y))\leq d(x,y)$. Ponadto
$d(H_u(y),H_v(y))=|u-v|d(\widetilde r,y)$.
Nierówność trójkąta daje więc
\[
d(H_u(x),H_v(y))\leq d(x,y)+|u-v|d(\widetilde r,y).
\]
Dowodzi to łącznej ciągłości względem punktu i parametru.
Mamy $H_0=\operatorname{id}$, $H_1\equiv\widetilde r$
i $H_u(\widetilde r)=\widetilde r$. Realizacja drzewa
jest zatem kontraktowalna, w szczególności jednospójna,
a skonstruowane nakrycie jest uniwersalne.
\end{proof}

\begin{uwaga}[Skończoność nakrycia]
Jeśli bazowy graf jest drzewem, jego nakrycie uniwersalne
jest z nim izomorficzne. Jeśli graf ma cykl, nakrycie
uniwersalne jest nieskończone. Wybierzmy cykl
(wliczając pętlę lub cykl z dwóch krawędzi równoległych)
i najkrótszą ścieżkę z $r$ do zbioru jego wierzchołków.
Jeżeli ścieżka jest niepusta, jej ostatnia krawędź
nie należy do cyklu. Po dojściu do cyklu przechodzimy
go w ustalonym kierunku dowolną liczbę razy.
Na styku kolejnych obiegów nie ma natychmiastowego
powrotu po tej samej krawędzi. Otrzymujemy różne
zredukowane spacery wychodzące z $r$ o nieograniczonych
długościach.
\end{uwaga}

\section{Nakrycia skończone i wzór Nielsena--Schreiera}
\label{sec:tg-finite-cover}

\begin{twierdzenie}[Charakterystyka Eulera nakrycia]
Niech $G$ będzie skończonym spójnym grafem, a
$p:|\widehat G|\to|G|$ spójnym, $d$-kartkowym nakryciem grafowym,
gdzie $d\ge1$ i $\widehat G$ jest grafem skończonym.
Wtedy $|V(\widehat G)|=d|V(G)|$ oraz
$|E(\widehat G)|=d|E(G)|$, a zatem
$\chi(|\widehat G|)=d\chi(|G|)$.
Jeśli $\pi_1(|G|)\cong F_r$, to
$\pi_1(|\widehat G|)\cong F_{1+d(r-1)}$.
\end{twierdzenie}

\begin{proof}
Każdy wierzchołek $v\in G$ ma dokładnie $d$ przeciwobrazów,
bo liczba kartek nakrycia jest stała na spójnej bazie.
Nad wnętrzem każdej krawędzi otrzymujemy dokładnie $d$
rozłącznych otwartych odcinków, stanowiących wnętrza
krawędzi $\widehat G$. To daje pierwsze dwie równości.
Odejmując liczby komórek wymiaru jeden od liczby
komórek wymiaru zero, dostajemy równość dla $\chi$.

Graf $\widehat G$ jest spójny, zatem jego liczba cyklomatyczna
wynosi
$|E(\widehat G)|-|V(\widehat G)|+1
=d(m-n)+1=1+d(r-1)$.
Z twierdzenia \ref{thm:tg-free-pi1} grupa fundamentalna
$\widehat G$ jest wolna o tej liczbie generatorów.
\end{proof}

\begin{przyklad}[Nakrycie trzykartkowe bukietu]
Niech $G$ będzie bukietem dwóch okręgów, przedstawionym jako
graf z jednym wierzchołkiem i dwiema pętlami. Dla dowolnego
spójnego nakrycia trzykartkowego $\widehat G$ mamy
$|V(\widehat G)|=3$, $|E(\widehat G)|=6$, a więc
$\pi_1(|\widehat G|)\cong F_{6-3+1}=F_4$.
Z twierdzenia o odpowiedniości spójnych nakryć
i podgrup grupy fundamentalnej wynika, że każda
podgrupa indeksu trzy grupy wolnej rzędu dwa jest
wolna rzędu cztery: bukiet jest spójny drogowo,
lokalnie spójny drogowo i półlokalnie jednospójny,
a nakrycie odpowiadające podgrupie dziedziczy
strukturę grafu przez podniesienie komórek.
\end{przyklad}

\begin{uwaga}[Subdywizje a nakrycia multigrafów]
Twierdzenie o grupie fundamentalnej obejmuje również
multigrafy i pętle. Można też sprowadzić je do przypadku
grafów prostych: wystarczy podzielić każdą krawędź w dwóch
różnych punktach jej wnętrza, otrzymując graf prosty
o tej samej realizacji topologicznej. Liczba $m-n+c$ nie zmienia się przy pojedynczej
subdywizji. Dlatego obliczenie dla bukietu z pętlami
jest szczególnym przypadkiem tej samej teorii.
\end{uwaga}

\begin{uwaga}[Nakrycia skończone drzew]
Jeśli $G$ jest drzewem, każde spójne nakrycie $|G|$
jest jedno-kartkowe. Rzeczywiście $|G|$ jest
po prostu spójna, więc każde spójne nakrycie
jest z nią izomorficzne. Dlatego przy $r=0$
wzór $1+d(r-1)$ z twierdzenia o charakterystyce Eulera nakrycia
nie daje ujemnych rang: dla $d>1$ spójne
nakrycie $d$-kartkowe drzewa w ogóle nie istnieje.
\end{uwaga}

\section{Opis nakryć skończonych za pomocą permutacji}
\label{sec:tg-permutation-covers}

Wzór na charakterystykę Eulera określa rangę grupy
fundamentalnej nakrycia, ale nie rozstrzyga, jak nakrycie
jest zbudowane. Poniższa konstrukcja pozwala opisać
również nakrycia, których nie można otrzymać przez
przesunięcia w grupie cyklicznej kartek.

\begin{twierdzenie}[Nakrycia i permutacje kartek]
\label{thm:tg-permutation-covers}
Niech $G=(V,E)$ będzie skończonym spójnym grafem
nieskierowanym (dopuszczamy pętle i krawędzie równoległe),
$T$ jego drzewem rozpinającym, a $d\geq1$ liczbą całkowitą.
Ustalmy orientację pomocniczą każdej krawędzi.
Każda rodzina permutacji $\sigma_e\in S_d$
dla $e\in E\setminus E(T)$ określa $d$-kartkowe
nakrycie grafowe przez następującą konstrukcję:
\begin{itemize}
\item wierzchołki to pary $(v,i)$, gdzie $v\in V$,
$i\in\{1,\ldots,d\}$;
\item dla każdej krawędzi $e:u\to v$ i każdego $i$
tworzymy osobną krawędź od $(u,i)$ do $(v,\sigma_e(i))$,
przyjmując $\sigma_e=\operatorname{id}$ dla $e\in E(T)$.
\end{itemize}
Każde $d$-kartkowe nakrycie grafowe $G$ jest
izomorficzne nad $G$ z nakryciem tej postaci.
Jest ono spójne dokładnie wtedy, gdy podgrupa
$K=\langle \sigma_e:e\notin E(T)\rangle\leq S_d$
działa przechodnio na $\{1,\ldots,d\}$.
Dwie rodziny $(\sigma_e)$ i $(\tau_e)$ wyznaczają
izomorficzne nakrycia nad $G$ dokładnie wtedy,
gdy istnieje jedna permutacja $\rho\in S_d$
taka, że $\tau_e=\rho\sigma_e\rho^{-1}$ dla każdego
$e\notin E(T)$.
\end{twierdzenie}

\begin{proof}
Rzutowanie na pierwszą współrzędną wierzchołka
i na nazwę krawędzi bazowej przedłużamy liniowo
na krawędziach realizacji. Nad każdym początkiem
$u$ krawędzi $e$ istnieje jeden jej początek
przy $(u,i)$, a nad końcem $v$ jeden koniec
przy $(v,j)$, ponieważ równanie $\sigma_e(i)=j$
ma dokładnie jedno rozwiązanie $i$.
Dla pętli liczymy te dwie incydencje oddzielnie,
nawet gdy tworzą tę samą pętlę w nakryciu.
Zatem rzutowanie spełnia lokalny warunek bijekcji
incydencji i jest surjektywne. Każde włókno nad
wierzchołkiem lub wnętrzem krawędzi ma $d$ punktów.

Odwrotnie, wybierzmy wierzchołek $r\in V$ i
ponumerujmy włókno nad $r$ liczbami $1,\ldots,d$.
Dla każdego $v\in V$ podnosimy jedyną ścieżkę w $T$
od $r$ do $v$ z każdego punktu włókna nad $r$.
Podnoszenie jest jednoznaczne: w każdym kroku
lokalna bijekcja wskazuje dokładnie jedną incydencję.
Odwrócenie ścieżki daje odwrotną bijekcję włókien,
więc końce podniesień numerują całe włókno nad $v$.
W tych numeracjach każda krawędź drzewa łączy
punkty o tym samym numerze; istotnie ścieżka do
jednego jej końca jest przedłużeniem ścieżki do
drugiego końca o tę krawędź.
Dla dowolnej pozostałej krawędzi $e$ lokalna
bijekcja wyznacza bijekcję włókien na jej końcach.
Jest to pewna permutacja $\sigma_e$ numerów.
Otrzymujemy dokładnie opisaną konstrukcję.

Każda ścieżka w nakryciu zmienia numer kartki
przez kolejne zastosowania $\sigma_e$ lub
$\sigma_e^{-1}$ (dla przejścia w kierunku przeciwnym).
Dlatego może połączyć tylko punkty o numerach
z tej samej orbity $K$. Odwrotnie, każdą permutację
$\sigma_e$ można zrealizować od włókna nad $r$
przez przejście w drzewie do początku $e$, przejście
$e$ i powrót w drzewie do $r$. Ścieżka odwrotna
realizuje $\sigma_e^{-1}$. Składając te ścieżki,
osiągamy wszystkie numery danej orbity.
Połączenia w drzewie pozwalają przejść następnie
do dowolnego wierzchołka bazy bez zmiany numeru.
Składowe spójne odpowiadają więc dokładnie orbitom $K$.

Wreszcie izomorfizm nad $G$ musi na włóknach
działać przez permutacje $\rho_v$.
Zgodność z krawędziami drzewa wymusza
$\rho_u=\rho_v$ dla każdej krawędzi $uv\in E(T)$,
więc wszystkie $\rho_v$ są jedną permutacją $\rho$.
Zgodność z krawędzią $e$ daje
$\rho(\sigma_e(i))=\tau_e(\rho(i))$ dla każdego $i$.
Jest to równoważne warunkowi sprzężenia w tezie.
Odwrotnie, gdy warunek zachodzi, odwzorowanie
$(v,i)\mapsto(v,\rho(i))$ wraz z odpowiednim
odwzorowaniem nazwanych krawędzi jest izomorfizmem
nad $G$.
\end{proof}

\begin{wniosek}[Przekształcenia pokładowe nakrycia]
\label{cor:tg-cover-centralizer}
W oznaczeniach twierdzenia \ref{thm:tg-permutation-covers}
grupa automorfizmów nakrycia nad $G$ jest izomorficzna
z centralizatorem
$C_{S_d}(K)=\{\rho\in S_d:\rho k=k\rho\ \text{dla każdego }k\in K\}$.
Jeżeli nakrycie jest spójne, to żaden nietrywialny
automorfizm nie ustala punktu włókna nad wierzchołkiem.
\end{wniosek}
\begin{proof}
W kryterium izomorfizmu przyjmujemy $\tau_e=\sigma_e$.
Warunek komutowania z generatorami jest równoważny
komutowaniu z całą grupą $K$, a składanie automorfizmów
odpowiada składaniu permutacji.
Jeśli $K$ działa przechodnio i $\rho(i)=i$, to dla
dowolnego $j$ istnieje $k\in K$ z $k(i)=j$.
Wówczas $\rho(j)=\rho(k(i))=k(\rho(i))=j$.
Zatem $\rho$ jest identycznością.
\end{proof}

\begin{przyklad}[Spójne nakrycie bez nietrywialnych przekształceń pokładowych]
\label{ex:tg-cover-nonregular}
Rozważmy bukiet dwóch okręgów z pętlami $a,b$ i trzy kartki.
Przyjmijmy $\sigma_a=(12)$ i $\sigma_b=(23)$.
Z kartki $1$ można przejść do $2$ przez $a$,
a następnie do $3$ przez $b$, więc działanie jest
przechodnie i nakrycie jest spójne. Permutacje te
generują $S_3$: oprócz identyczności i dwóch
transpozycji otrzymujemy $(13)$ oraz oba cykle
długości trzy.
Permutacja komutująca z $(12)$ musi zachowywać
jej zbiór punktów stałych $\{3\}$, a komutująca
z $(23)$ musi ustalać $1$. Musi zatem ustalać
również $2$. Centralizator jest trywialny,
mimo że nakrycie ma trzy kartki.
Z wcześniejszego wzoru jego grupa fundamentalna
jest izomorficzna z $F_4$.
\end{przyklad}

\section{Metryka długości na grafie}
\label{sec:tg-metric}

\begin{definicja}[Metryka długości]
Niech $G$ będzie spójnym grafem skończonym. Nadajmy każdej
krawędzi długość jeden i utożsammy ją z odcinkiem $[0,1]$.
Długość krzywej w $|G|$ liczymy odcinkami euklidesowymi;
przez $d_{|G|}(x,y)$ oznaczamy infimum długości krzywych
łączących $x,y\in|G|$.
\end{definicja}

\begin{stwierdzenie}[Zgodność metryki grafowej z metryką długości]
Dla spójnego skończonego grafu $G$ i dowolnych wierzchołków
$u,v\in V(G)$ zachodzi $d_{|G|}(u,v)=d_G(u,v)$.
\end{stwierdzenie}

\begin{proof}
Kombinatoryczna najkrótsza ścieżka długości $k=d_G(u,v)$
po sklejeniu odcinków jest krzywą w $|G|$ długości $k$,
a zatem $d_{|G|}(u,v)\le k$.

Dla nierówności odwrotnej przedłużmy liniowo na krawędziach
funkcję $f(w)=d_G(u,w)$ z wierzchołków. Nierówność trójkąta
daje $|f(a)-f(b)|\le1$ na końcach krawędzi $ab$, więc
$f$ jest $1$-Lipschitzowska na każdej krawędzi.
W małej gwieździe wierzchołka zmiana $f$ między ramionami
jest ograniczona sumą odległości od wierzchołka.
Dzieląc parametr krzywej $\gamma$ na skończenie wiele
przedziałów o obrazach w takich gwiazdach lub wnętrzach
krawędzi (z lematu o liczbie Lebesgue'a), otrzymujemy
$k=|f(v)-f(u)|\le\operatorname{len}(\gamma)$.
Infimum po krzywych daje $d_{|G|}(u,v)\ge k$.
\end{proof}

